\section{讨论与后续检验计划}
两项实验中的一项关键对照，都使用了这样两条历史：观测相同，要求的答案却不同。如果系统只能看到两条历史共有的部分，就不可能在两者上都答对。利用这样的历史对，可以确定哪些信息是必需的。要进一步断言某个具体系统确实使用了这些信息，还须在固定后续输入时检验传递、替换或移除载体的效果。单凭对依赖锥的几何覆盖、可从数据中恢复的相关量或较高的平均准确率，都不足以确立这样的效果。

\paragraph{结论的适用范围。}
无差错编码与图着色之间的关系是信息论中的已知结果 \cite{witsenhausen,charpenay}。本文的贡献在于设计可检验的干预、给出具体的离散见证，并限制两个回合之间的通道。在无限直线上，对每个时刻 $t$ 和窗口内的每个坐标 $j$，都存在一对仅在 $j$ 处有别的配置，使目标比特不同。若 $j$ 在窗口外，对于任意初始配置，仅改变该处的值都会保持目标比特不变。在 OLMo 实验中，表内每一行都对应特定的冻结检查点（Base 或 Instruct）与接口的组合，结果只涉及同一个有限测试集。判定系统具有或缺乏主观体验，都超出了这些数据所能支持的范围。

\paragraph{向其他领域的迁移。}
在机制审计中，应先检验模型在历史完整时能否完成任务，然后可以转移候选的隐藏载体，并检验擦除它会有何影响。在意识研究中，包括对意识障碍的讨论，宜分别考察报告能否取得、关于体验作了哪些假设，以及是否采取预防措施。在临床上应用这种区分仍需独立验证。状态的来源也可成为数字取证与分布式计算中的研究问题。面对固定查询，正确的预录记录可以复现实时应答。在一组改变地址的独立试验中，实时通道答对 8/8，预录记录答对 6/8。分支的来历与其能否回答新查询可以分别检验。

\paragraph{OLMo 中的可寻址反馈（T3，\texttt{NOT\_RUN}）。}
在完成第 33--44 名组成的测试集后，本实验需要另设测试集。模型保留此前的两个选择 $h_0,h_1$。它须从新任务的内容推断哪一个选择需要修正，标准答案中的地址 $K$ 对模型隐藏。在看到输入串 $\mathbf{x}$ 前，模型自行给出一个查询比特 $q$。观察者保存着两个反馈比特 $e_0,e_1$。它返回其中一个，随后模型须对输入串应用目标规则 $h_K\oplus e_K$。模型获知反馈比特的唯一途径是观察者的回答。

对于相同的历史，向观察者传递地址 $q$（实时查询）、$0$（固定地址）或 $1-q$（地址替换）。接收方在所有条件下均在同一位置看到一个比特。实际传递的地址和标记替换与否的信息均记入日志，并对接收方隐藏。分别测量是否选中有用的查询（$q=K$）以及 $L\ne R$ 时最终答案的准确率。在 $e_0\ne e_1$ 时作配对比较，可以检验答案是否随地址替换而改变。

在端点不同的输入串上，$K$ 的两个取值出现次数相同。对每个任务、每一对过去的选择和每个输入串，四种反馈比特组合也均等出现。固定地址时准确率的上限为 $75\%$。当 $K=0$ 时，返回的 $e_0$ 正是所需反馈。当 $K=1$ 时，即使收到了 $e_0$，$e_1$ 的两个取值仍然等概率。因此最优准确率为 $\tfrac12\cdot1+\tfrac12\cdot\tfrac12=\tfrac34$。应在由端点不同的输入串组成的完整平衡测试集上，检验实时查询能否超过这个上限。

\begin{center}\small
\begin{tabularx}{\linewidth}{@{}p{.23\linewidth}X@{}}\toprule
本版本的任务 & 所需证据 \\\midrule
G4-1，第 33--44 名 & 在首次启用共同留出集前，在开发集上固定待比较的两个接口。精确的一比特转移要求 12/12 个整块全部通过，并通过擦除、格式和通道边界对照。\\
G3，第 21--32 名 & 文本通道通过检验后，用一次预先规定的激活脉冲替代文本。关闭脉冲时，在相同选择和输入下，两条反馈分支须一致。\\
AI.04：多维 G4 与 DAS & 在原始位置不可访问时转移指定的内部状态。要定位子空间，须预先指定干预层、中介度量、容差和独立对照。检验 $C_{\rm med}$ 还需要使用同一规则的另一种任务。\\
R110.11，条目 A6 & 在同一空间均匀的二值 Rule~110 系统内实现存储和回送。已证实的自主构造使用扩展字母表。\\
合并与回送的形式化 & 在 Lean 中形式化对 4096 个环境的穷举以及从 $t=4\to5$ 的构造性见证。几何证明的根定理和观察者模块 11/11 项查询的验证已复现（附录~\ref{app:lean}）。\\
$D_{\rm dyn}$ & 证明组织结构的变化能在回合结束后持续，并在受控重置下展示该变化对新输入的作用。\\
\bottomrule\end{tabularx}
\end{center}
附录~\ref{app:gates}~中的规约列出了所提出的通道检验失败的条件。某一通道的阴性结果限制了对所选载体可以作出的推断。相关假设及其状态见附录~\ref{app:registers}。
